\begin{defi}\label{Defn: weight function}
Consider a set of variables $\bm w\coloneqq (w_{\mathcal{C}})_{(\mathcal{C}\in\mathcal{A}/W)}$ and a weight function $$\op{wt}:\mathcal{R}\rightarrow\{ w_{\mathcal{C}}\ |\ \mathcal{C}\in\mathcal{A}/W\},$$ such that $\op{wt}(t)=w_{\mathcal{C}}$ if $\mathcal{C}$ is the orbit that contains the fixed hyperplane $V^t$. Then, the weighted enumeration of reflection factorizations of some element $g\in W$ is encoded via the following generating function:
$$\op{FAC}_{W,g}(\bm w,z)\coloneqq \sum_{\substack{(t_1,\cdots,t_N)\in\mathcal{R}^N\\t_1\cdots t_N=g}}\op{wt}(t_1)\cdots\op{wt}(t_N)\cdot \dfrac{z^N}{N!}. $$
\end{defi}

Because the sets $\mathcal{C}^{\op{ref}}\coloneqq \{t\in\mathcal{R}\ |\ V^t\in\mathcal{C}\}$ are closed under conjugation, the Lemma of Frobenius can again be used to express $\op{FAC}_{W,g}(\bm w,z)$ as a finite sum of exponentials. Notice first, that the order of the subsets $A_i$ in Thm.~\ref{Thm: Frobenius lemma} does not affect the enumeration as the different sets of factorizations have the same size. Indeed, one can easily construct a bijective map by considering a sequence of {\it Hurwitz moves}:$$
(t_1,t_2,\ldots,t_k,t_{k+1},\ldots ,t_l)\rightarrow (t_1,t_2,\ldots,t_kt_{k+1}t_k^{-1},t_k,\ldots,t_l).$$  

Having said that, and assuming there are $r=|\mathcal{A}/W|$ different orbits of hyperplanes, denoted $\mathcal{C}_1,\ldots, \mathcal{C}_r$, Thm.~\ref{Thm: Frobenius lemma} now implies that \begin{align*}
\op{FAC}_{W,g}(\bm w,z)=&\sum_{\substack{ N\geq 0\\l_1+\cdots +l_r=N}}\binom{N}{l_1,\ldots,l_r}w_{\mathcal{C}_1}^{l_1}\cdots w_{\mathcal{C}_r}^{l_r} \dfrac{z^N}{N!}\times\\
&\times\dfrac{1}{|W|}\sum_{\chi\in\widehat{W}}\chi(1)\cdot \chi(g^{-1})\cdot \Big[ \dfrac{\chi(\mathcal{C}^{\op{ref}}_1)}{\chi(1)}\Big]^{l_1}\cdots \Big[ \dfrac{\chi(\mathcal{C}^{\op{ref}}_r)}{\chi(1)}\Big]^{l_r}.\end{align*}Using standard properties of exponential generating functions, we can rewrite the sum as$$ 
\op{FAC}_{W,g}(\bm w,z)=\dfrac{1}{|W|}\sum_{\chi\in\widehat{W}}\chi(1)\cdot\chi(g^{-1})\cdot\op{exp}\Big[ zw_{\mathcal{C}_1}^{\phantom{l_1}}\cdot \dfrac{\chi(\mathcal{C}_1^{\op{ref}})}{\chi(1)}\Big]\cdots \op{exp}\Big[ zw_{\mathcal{C}_r}^{\phantom{l_r}}\cdot \dfrac{\chi(\mathcal{C}_r^{\op{ref}})}{\chi(1)}\Big].$$
Finally, notice that by Defn.~\ref{Defn: Local Coxeter numbers} we can rewrite the quantities in the exponentials in terms of local Coxeter numbers. Indeed, we have $c_{\chi_{,\mathcal{C}}}=|\mathcal{C}^{\op{ref}}|-\chi(\mathcal{C}^{\op{ref}})/\chi(1)$ and if we define $\op{wt}(\mathcal{R})\coloneqq \sum_{t\in\mathcal{R}}\op{wt}(t)$, the previous expression becomes a direct analog of \eqref{EQ: Frobenius via Coxeter numbers}: 
\begin{equation}
\op{FAC}_{W,g}(\bm w,z)=\dfrac{e^{z\cdot\op{wt}(\mathcal{R})}}{|W|}\sum_{\chi\in\widehat{W}}\chi(1)\cdot\chi(g^{-1})\cdot \big(e^{-zw_{\mathcal{C}_1}^{\phantom{l_1}}}\big)^{\displaystyle c_{\chi_{,\mathcal{C}_1}}}\cdots \big( e^{-zw_{\mathcal{C}_r}^{\phantom{l_r}}}\big)^{\displaystyle c_{\chi_{,\mathcal{C}_r}}}.\label{EQ: weighted enum via local Coxeter numbers}
\end{equation}



\begin{lemma}\label{Lemma: weigted enumeration only multiples of |g| contribute}
For a complex reflection group $W$, and a \textbf{regular} element $g\in W$, the total contribution in \eqref{EQ: weighted enum via local Coxeter numbers} of those characters $\chi\in\widehat{W}$ for which \textbf{any} $c_{\chi_{,\mathcal{C}}}$ is not a multiple of $|g|$ is $0$.
\end{lemma}
\begin{proof}
The proof is essentially the same as for Lemma~\ref{Lemma: only multiples of |g| contribute}. However, we first need to order the orbits $\mathcal{C}\in\mathcal{A}/W$ (arbitrarily) and then apply the same idea sequentially. 

We start by partitioning the set of irreducible characters $\chi\in\widehat{W}$ into orbits under the action of $\Psi_{\mathcal{C}_1}$. By Prop.~\ref{Prop: properties of Psi for Cox numbers} all characters in such an orbit have the same local Coxeter numbers. Pick a character $\chi$ whose local Coxeter number $c_{\chi_{,\mathcal{C}_1}}$ is not a multiple of $|g|$; by Prop.~\ref{Prop: The key technical lemma} we may further assume that $\chi(g)\neq 0$. Let $k$ be the smallest number such that $\Psi_{\mathcal{C}_1}^k(\chi)=\chi$; again we will have by Prop.~\ref{Prop: The key technical lemma} that $k$ must be a multiple of $m\coloneqq \frac{|g|}{\op{gcd}(c_{\chi_{,\mathcal{C}_1}},|g|)}$. Now, since by Prop.~\ref{Prop: properties of Psi for Cox numbers} the degrees of characters and the local Coxeter numbers are respected by $\Psi_{\mathcal{C}_1}$, it is enough to show that $$\sum_{j=1}^k\Psi^j_{\mathcal{C}_1}(\chi)(g^{-1})=0.$$ 
Indeed, this follows immediately from Prop.~\ref{Prop: The key technical lemma} as $\Psi_{\mathcal{C}_1}^j(\chi)(g^{-1})=\xi^j\chi(g^{-1})$ for the primitive $m$-th root of unity $\xi\coloneqq \zeta^{-c_{\chi_{,\mathcal{C}_1}}}$ (and since $m\neq 1$ divides $k$). Notice now that we can continue like this, eventually disregarding all characters with local Coxeter number $c_{\chi_{,\mathcal{C}_1}}$ not a multiple of $|g|$. Then we may proceed with the {\it remaining} characters and the orbit $\mathcal{C}_2$ and since, by Prop.~\ref{Prop: properties of Psi for Cox numbers}, $\Psi_{\mathcal{C}_2}$ also respects the numbers $c_{\chi_{,\mathcal{C}_1}}$ (which now \emph{for all remaining characters} are multiples of $|g|$) we do not have to worry that we might eventually cancel the same character twice. We go on like this with all orbits $\mathcal{C}_i$ and the proof is complete.
\end{proof}
